% GENERATED FILE. Edit canonical lessons, then run npm run generate:books.
% canonical-source-hash: fnv1a64:f95aa21f5e8753b8
% canonical-lessons: TE-C137-telika, TE-C137-mettani, TE-C137-gatti, TE-C137-subhram, TE-C137-muriki

\chapter{Light, Soft, Hard, Clean, Dirty}
\label{ch:te-light-soft-hard-clean-dirty}
\hlchaptermodality{137}

\begin{chapteropening}
\textbf{By the end of this chapter:} \emph{I can say light, soft, hard, clean and dirty.}

Complete the last lesson of chapter 137: i can say light, soft, hard, clean and dirty.
\end{chapteropening}

\section[\texorpdfstring{\te{తేలిక} (tēlika)}{తేలిక (tēlika)}]{\te{తేలిక} (tēlika) — light (in weight)}

\label{lesson:TE-C137-telika}



\begin{quote}
Before the new one: say the Telugu for short, then the Telugu for heavy.
\end{quote}

\subsection*{You'll want to know: \te{తేలిక}}
\textbf{\te{తేలిక}} — \emph{tēlika} — \textquotedblleft{}light (in weight)\textquotedblright{}.

\begin{grammarlens}[title={What you've built}]
One more word for this chapter.
\end{grammarlens}

\subsection*{Your turn}
\begin{itemize}
\raggedright
  \item \emph{Say it:} \emph{tēlika}
  \item \emph{Say it:} \emph{tēlika}, once more
  \item \emph{Say it:} say \emph{tēlika}
  \item \emph{Recall:} say the Telugu for short, then the Telugu for heavy, then say \emph{tēlika} again
\end{itemize}

\begin{tcolorbox}[breakable,colback=teal!4,colframe=teal!35!black,title={Before you move on}]
What does \te{తేలిక} mean? (\textquotedblleft{}Light (in weight)\textquotedblright{}.) Say it once more.
\end{tcolorbox}
\section[\texorpdfstring{\te{మెత్తని} (mettani)}{మెత్తని (mettani)}]{\te{మెత్తని} (mettani) — soft}

\label{lesson:TE-C137-mettani}



\begin{quote}
Before the new one: say the Telugu for heavy, then the Telugu for light.
\end{quote}

\subsection*{You'll want to know: \te{మెత్తని}}
\textbf{\te{మెత్తని}} — \emph{mettani} — \textquotedblleft{}soft\textquotedblright{}.

\begin{grammarlens}[title={What you've built}]
One more word for this chapter.
\end{grammarlens}

\subsection*{Your turn}
\begin{itemize}
\raggedright
  \item \emph{Say it:} \emph{mettani}
  \item \emph{Say it:} \emph{mettani}, once more
  \item \emph{Say it:} say \emph{mettani}
  \item \emph{Recall:} say the Telugu for heavy, then the Telugu for light, then say \emph{mettani} again
\end{itemize}

\begin{tcolorbox}[breakable,colback=teal!4,colframe=teal!35!black,title={Before you move on}]
What does \te{మెత్తని} mean? (\textquotedblleft{}Soft\textquotedblright{}.) Say it once more.
\end{tcolorbox}
\section[\texorpdfstring{\te{గట్టి} (gaṭṭi)}{గట్టి (gaṭṭi)}]{\te{గట్టి} (gaṭṭi) — hard}

\label{lesson:TE-C137-gatti}



\begin{quote}
Before the new one: say the Telugu for light, then the Telugu for soft.
\end{quote}

\subsection*{You'll want to know: \te{గట్టి}}
\textbf{\te{గట్టి}} — \emph{gaṭṭi} — \textquotedblleft{}hard\textquotedblright{}.

\begin{grammarlens}[title={What you've built}]
One more word for this chapter.
\end{grammarlens}

\subsection*{Your turn}
\begin{itemize}
\raggedright
  \item \emph{Say it:} \emph{gaṭṭi}
  \item \emph{Say it:} \emph{gaṭṭi}, once more
  \item \emph{Say it:} say \emph{gaṭṭi}
  \item \emph{Recall:} say the Telugu for light, then the Telugu for soft, then say \emph{gaṭṭi} again
\end{itemize}

\begin{tcolorbox}[breakable,colback=teal!4,colframe=teal!35!black,title={Before you move on}]
What does \te{గట్టి} mean? (\textquotedblleft{}Hard\textquotedblright{}.) Say it once more.
\end{tcolorbox}
\section[\texorpdfstring{\te{శుభ్రం} (śubhraṁ)}{శుభ్రం (śubhraṁ)}]{\te{శుభ్రం} (śubhraṁ) — clean}

\label{lesson:TE-C137-subhram}



\begin{quote}
Before the new one: say the Telugu for soft, then the Telugu for hard.
\end{quote}

\subsection*{You'll want to know: \te{శుభ్రం}}
\textbf{\te{శుభ్రం}} — \emph{śubhraṁ} — \textquotedblleft{}clean\textquotedblright{}.

\begin{grammarlens}[title={What you've built}]
One more word for this chapter.
\end{grammarlens}

\subsection*{Your turn}
\begin{itemize}
\raggedright
  \item \emph{Say it:} \emph{śubhraṁ}
  \item \emph{Say it:} \emph{śubhraṁ}, once more
  \item \emph{Say it:} say \emph{śubhraṁ}
  \item \emph{Recall:} say the Telugu for soft, then the Telugu for hard, then say \emph{śubhraṁ} again
\end{itemize}

\begin{tcolorbox}[breakable,colback=teal!4,colframe=teal!35!black,title={Before you move on}]
What does \te{శుభ్రం} mean? (\textquotedblleft{}Clean\textquotedblright{}.) Say it once more.
\end{tcolorbox}
\section[\texorpdfstring{\te{మురికి} (murikī)}{మురికి (murikī)}]{\te{మురికి} (murikī) — dirty}

\label{lesson:TE-C137-muriki}



\begin{quote}
Before the new one: say the Telugu for hard, then the Telugu for clean.
\end{quote}

\subsection*{You'll want to know: \te{మురికి}}
\textbf{\te{మురికి}} — \emph{murikī} — \textquotedblleft{}dirty\textquotedblright{}.

\begin{grammarlens}[title={What you've built}]
One more word for this chapter.
\end{grammarlens}

\subsection*{Your turn}
\begin{itemize}
\raggedright
  \item \emph{Say it:} \emph{murikī}
  \item \emph{Say it:} \emph{murikī}, once more
  \item \emph{Say it:} say \emph{murikī}
  \item \emph{Recall:} say the Telugu for hard, then the Telugu for clean, then say \emph{murikī} again
\end{itemize}

\begin{tcolorbox}[breakable,colback=teal!4,colframe=teal!35!black,title={Before you move on}]
What does \te{మురికి} mean? (\textquotedblleft{}Dirty\textquotedblright{}.) Say it once more.
\end{tcolorbox}
